\section{Simulations and Performance Modeling}

Simulation and modeling played a central role in validating the feasibility of the thrust vector control system prior to experimental flight testing. Given the risks associated with uncontrolled instability in TVC rockets, analytical and numerical models were used to predict vehicle behavior under both ideal and disturbed conditions. These simulations served to verify that the selected motor, gimbal geometry, servo dynamics, and control logic could theoretically maintain stability throughout flight, and to establish performance expectations against which experimental data could later be compared. Rather than relying on a single simulation framework, a hybrid modeling approach was used: OpenRocket \cite{openrocket2023} was used to analyze baseline trajectory behavior and aerodynamic stability using established rocketry methods, while custom analytical models were developed to evaluate thrust vector control authority, angular response, and recovery limits. This separation allowed aerodynamic and control effects to be studied independently before being considered as a coupled system.



\subsection{Baseline Trajectory and Aerodynamic Stability Modeling}

The first stage of simulation focused on establishing a baseline flight profile for the rocket in the absence of active control. This reference case provided expected values for apogee, velocity, and stability margin, and served as a benchmark for evaluating the effectiveness of thrust vector control.
OpenRocket was selected for this task due to its widespread acceptance in amateur and educational rocketry and its use of Barrowman-based aerodynamic modeling \cite{barrowman1967}. The rocket geometry, mass distribution, and propulsion characteristics described in Section 2 were replicated within the software environment. Particular care was taken to ensure that the simulated center of mass matched the experimentally measured value and that the thrust curve corresponded to the manufacturer-provided data for the Estes E12-4 motor. The result of this simulation confirmed that, in the absence of fins or active control, the vehicle was aerodynamically unstable: the center of pressure remained ahead of the center of mass under static conditions, producing a negative static margin. Without restoring aerodynamic surfaces, this configuration caused even minor disturbances to grow rather than decay, consistent with the divergent, inverted-pendulum-like behavior described in Section 2.5. This outcome reinforced the necessity of active stabilization and justified the exclusive reliance on thrust vector control for attitude correction.


\subsection{Thrust Vector Control Authority Modeling}

To determine whether the thrust vectoring system could realistically stabilize the rocket during powered flight, an analytical model of control authority was developed. Control authority was defined as the maximum angular acceleration the TVC system could impart on the vehicle, given the available thrust (derived from the thrust curve) and the limitations of gimbal geometry. 
The corrective moment generated by thrust vectoring was modeled as a function of instantaneous motor thrust, the perpendicular offset between the thrust vector and the center of mass, and the gimbal deflection angle. Using the relationship introduced in Sections 3.3 and 3.4.1, the maximum available corrective torque was calculated across the duration of the motor burn using the thrust curve. This analysis revealed that control authority was highest shortly after ignition, when thrust peaked, and decreased steadily as the motor approached burnout. 
This time-varying nature of control authority imposed an important constraint on recovery capability. While the system could correct relatively large disturbances early in the burn, its ability to counteract angular deviations diminished rapidly as thrust decreased. As a result, the control system was most effective during the early powered phase of flight, when both thrust and corrective torque were highest. This observation directly informed the design of the launch detection and apogee logic described in Section 4, ensuring that control efforts were concentrated during the period of maximum effectiveness.


\subsection{Disturbance Modeling and Tip-Off Response}

In addition to idealized vertical flight, the rocket’s response to common disturbances was evaluated through simplified dynamic models. The most critical of these disturbances was launch rail tip-off, which occurs when the rocket exits the launch guide with a small angular deviation and nonzero angular velocity. Due to the absence of fins, even modest tip-off angles can grow rapidly if not corrected early in flight.

Angular dynamics were approximated using a rigid-body rotational model, where the angular acceleration was determined by the net torque acting on the vehicle and its moment of inertia. Disturbance torques arising from aerodynamic forces and gravity were compared against the available corrective torque from thrust vectoring. These simulations demonstrated that for small initial tilt angles, the PID-controlled TVC system could arrest angular growth and return the rocket toward vertical alignment within a fraction of a second. However, beyond a critical initial tilt, angular velocity increased faster than the system could counteract, leading to divergence consistent with the recovery limits discussed in Section 3.4.

These results emphasized the importance of early control engagement and justified the conservative launch detection thresholds implemented in the flight software. By enabling thrust vector control only after sustained acceleration and stable sensor readings were detected, the system minimized the risk of reacting too late to rapidly growing disturbances.


\subsection{Simulation-Informed Control Parameter Selection}

The results of the trajectory and disturbance simulations were used to constrain the range of acceptable PID gains and gimbal deflection limits. Simulations showed that aggressive proportional gains could, in theory, reduce recovery time but often resulted in actuator saturation and oscillatory behavior due to servo response limitations. Conversely, overly conservative gains improved stability but reduced the system’s ability to counteract rapid disturbances during launch. Therefore, it was vital to select PID gains which would find a balance between these two extremes.

As a result, PID parameters were selected to balance responsiveness and stability, prioritizing predictable recovery behavior over maximum correction speed. Integral gain was limited to prevent windup during periods of actuator saturation, while derivative gain was tuned to damp oscillations without amplifying high-frequency sensor noise. These choices ensured that the simulated system behavior aligned with the physical constraints of the servos and the thrust vectoring mechanism.


\subsection{Role of Simulation in Experimental Validation}

Ultimately, the simulations developed in this section were not intended to perfectly predict flight behavior, but rather to define a realistic operating envelope for the thrust vector control system. By identifying control authority limits, disturbance sensitivity, and recovery boundaries in advance, the simulations reduced the risk of catastrophic failure during experimental testing.


